\section{5. Running exploration}

\subsection{The exploration cycle}

\begin{longtable}{p{0.34\linewidth}p{0.6\linewidth}}

\rowcolor{tablehead} \textbf{Sequence} & \textbf{Resolution} \\

\endfirsthead

\rowcolor{tablehead} \textbf{Sequence} & \textbf{Resolution} \\

\endhead

1 & Explore the current chamber \\

2 & Inspect and choose a door \\

3 & Resolve locks and applicable traps \\

4 & Open door and shift dungeon \\

5 & Describe newly revealed chamber \\

6 & Players decide whether to cross \\

7 & Close door \\

8 & Resolve collapse, then pending doomsday \\

9 & Continue exploration \\

\end{longtable}

Avoid treating every room transition as an extended mechanical procedure. Once players understand the process, resolve routine actions quickly. Reserve detailed narration and initiative tracking for situations where decisions carry meaningful consequences.

\subsection{Describing rooms}

When a door opens, describe the chamber beyond using its player-facing appearance text. The room is discovered at this moment, even if nobody enters. Describe visible architecture, inhabitants, objects, sounds, smells and obvious hazards. Do not automatically reveal hidden mechanisms, puzzle solutions, enemy weaknesses or room classification. Separate DM descriptions should not be read aloud.

\subsection{Mapping the dungeon}

Allow players to draw maps, record room names, mark doors and track previous movements. Do not discourage mapping simply because the dungeon shifts. Understanding how and where chambers move is an important strategic skill. Keep the DM application private: it contains information about undiscovered chambers, hidden traps and treasure. A player-generated map may become inaccurate after a shift; this is intended.

\subsection{Splitting the party}

Characters may occupy different rooms. Splitting can provide tactical advantages but introduces substantial risks. A character may become separated from companions when surrounding chambers shift. Communication and regrouping may become difficult, especially during doomsday. Alternate between groups frequently and use a shared initiative order during simultaneous encounters if necessary.
